\section{The queue is an interval, not a pattern of zero bytes}\label{maze-p05}

At a checkpoint, the queue's active work lies between head, inclusive, and tail, exclusive. The prefix before head is retained processing history. The suffix at and beyond tail is unused storage. A zero byte in either region does not identify an empty slot, because zero is a valid cell ID.

The counters increase and do not wrap. That choice is safe because a cell becomes waiting when first discovered and never returns to unseen. At most 64 discovery operations can append entries. Thus head and tail remain between zero and 64, and an append uses an index below 64. Tail can equal 64 once the final available slot has been filled.

Retaining the prefix makes the teaching trace easy to inspect and lets the grader verify FIFO order. This queue does not yet solve the general problem of reusing bounded storage for an indefinitely long stream. The ring-buffer assignment later in the book addresses that different lifetime.
